%! Mode:: "TeX:UTF-8"
%! TEX program = xelatex
\documentclass[withoutpreface,bwprint]{cumcmthesis}

\usepackage{etoolbox}
\usepackage[numbers,sort&compress]{natbib}
\usepackage{url}
\usepackage{booktabs}
\usepackage{tabularx}
\usepackage{array}
\usepackage{caption}
\usepackage{fancyhdr}
\usepackage{enumitem}
\BeforeBeginEnvironment{tabular}{\zihao{-5}}
\BeforeBeginEnvironment{tabularx}{\zihao{-5}}

% A4 纸，四边 2.5 cm；正文小四号、1.25 倍行距，中西文字体分别为宋体和 Times New Roman。
\geometry{a4paper,top=2.5cm,bottom=2.5cm,left=2.5cm,right=2.5cm}
\renewcommand{\baselinestretch}{1.25}
\setlength{\parindent}{2em}
\setmainfont{Times New Roman}
\setCJKmainfont{SimSun}

% 无页眉；页码位于页脚中部。标题/摘要页不显示页码，正文从第二页显示 1。
\pagestyle{fancy}
\fancyhf{}
\fancyfoot[C]{\thepage}
\renewcommand{\headrulewidth}{0pt}
\renewcommand{\footrulewidth}{0pt}
\fancypagestyle{plain}{\fancyhf{}\fancyfoot[C]{\thepage}\renewcommand{\headrulewidth}{0pt}\renewcommand{\footrulewidth}{0pt}}

% 题目三号黑体；一级标题四号黑体居中；二、三级标题小四号黑体左对齐。
\titleformat{\section}{\centering\heiti\zihao{4}}{\thesection}{1em}{}
\titleformat{\subsection}{\heiti\zihao{-4}}{\thesubsection}{1em}{}
\titleformat{\subsubsection}{\heiti\zihao{-4}}{\thesubsubsection}{1em}{}
\DeclareCaptionFont{wuh}{\heiti\zihao{-5}}
\captionsetup{font=wuh,labelfont=wuh,justification=centering}
\captionsetup[figure]{font=wuh,labelfont=wuh,justification=centering}
\bibliographystyle{gbt7714-numerical}

\title{论文题目（请填写）}

\begin{document}
\pagenumbering{arabic}
\setcounter{page}{0}
\maketitle
\thispagestyle{empty}

\begin{abstract}
请在此填写中文摘要。摘要应简明说明研究问题、建模思路、主要方法、核心结果与结论，避免出现单位、姓名、队号等可能泄露答题人身份的信息。

\noindent\textbf{关键词：}关键词一；关键词二；关键词三；关键词四
\end{abstract}

% 正文从此页开始，页脚中部显示阿拉伯数字 1。
\setcounter{page}{1}

\section{问题重述}
简要、客观地重述题目要求，不照抄题面；涉及公开资料时，在正文中按顺序编码引用，例如文献 \cite{example}。

\section{问题分析}
说明问题的结构、关键难点、总体思路以及各子问题之间的关系。

\section{模型假设}
在此列出合理、必要且可检验的模型假设。

\section{符号说明}
\begin{table}[htbp]
  \centering
  \caption{主要符号说明}
  \begin{tabularx}{0.9\textwidth}{>{\centering\arraybackslash}p{0.22\textwidth}X}
    \toprule
    符号 & 含义 \\
    \midrule
    $x$ & 请填写符号含义 \\
    $y$ & 请填写符号含义 \\
    \bottomrule
  \end{tabularx}
\end{table}

\section{模型建立与求解}
\subsection{子问题一模型}
请填写模型建立、公式推导、算法设计和求解过程。

\subsection{子问题二模型}
请填写模型建立、公式推导、算法设计和求解过程。

\subsubsection{算法流程}
请插入算法流程图或伪代码。图题、表题使用五号黑体并居中，图、表编号分别连续编号。

\section{结果分析与模型检验}
展示计算结果、敏感性分析、误差分析、稳定性检验或对比实验。

\section{模型评价与推广}
\subsection{模型优点}
请填写。
\subsection{模型不足}
请填写。
\subsection{模型推广}
请填写。

\section{结论}
请用简洁文字概括主要结论，确保结论均能由前文模型和结果支持。

\section*{AI工具使用情况说明}
请在此说明论文写作、程序编写、数据处理或排版过程中使用的 AI 工具及其具体用途。

% 参考文献置于结论之后、附录之前。
\bibliography{ref}

% 附录标题与第一张图同页，图片不加外框，四张题干图片各占一页。
\clearpage
\section*{附录}

\end{document}
